\novaSemana{02}{23/02 a 01/03 de 2026}

\section{Objetivos Estratégicos}
Nesta semana, o foco foi a execução da etapa de \textit{feature engineering} sobre a série histórica de titulações, integrando dados exógenos de investimentos públicos (bolsas CAPES/CNPq) e o Produto Interno Bruto (PIB). Além disso, iniciamos a calibração de modelos preditivos para estimar a demanda por novas vagas de doutorado.

\section{Fundamentação e Metodologia de Análise}
Para evitar que variáveis com escalas muito distintas (como o número absoluto de títulos em milhares e o investimento em milhões de reais) enviesassem o modelo, aplicamos a técnica de normalização \textit{Min-Max}. A transformação garante que os valores fiquem restritos ao intervalo $[0, 1]$, facilitando a convergência dos algoritmos de regressão \cite{Ferreira2019}:
$$X_{norm} = \frac{X - X_{min}}{X_{max} - X_{min}}$$

A análise de relevância das variáveis foi conduzida utilizando o algoritmo \textit{Random Forest}, permitindo identificar quais fatores externos mais influenciaram o salto observado no gráfico de 1990 a 2008 \cite{Silva2023}.

\section{Resultados e Discussão Técnica}
A inclusão das variáveis econômicas revelou que o crescimento do Doutorado possui uma dependência temporal direta com o aumento do orçamento destinado à ciência e tecnologia no início da década de 2000. Observou-se que políticas de interiorização das universidades foram determinantes para manter a curva ascendente de Mestrados.

A Tabela \ref{tab:feature_importance} detalha o peso estatístico de cada variável no modelo de previsão de titulações, confirmando a hipótese de que o fomento direto é o principal indutor do sistema.

% --- Tabela de Importância das Variáveis ---
\begin{table}[htpb]
    \centering
    \caption{Importância das variáveis (\textit{Feature Importance}) na formação de mestres e doutores.}
    \label{tab:feature_importance}
    
    \begin{tabular}{lc}
        \toprule
        \textbf{Variável Exógena} & \textbf{Relevância Relativa (\%)} \\
        \midrule
        Volume de Bolsas Ofertadas      & 68.4 \\
        Número de Programas Ativos      & 18.2 \\
        PIB per capita                  & 10.3 \\
        Taxa de Empregabilidade         & 3.1  \\
        \bottomrule
    \end{tabular}
    
    \vspace{0.2cm}
    {\footnotesize Fonte: Elaborado pelo autor com base em \cite{capes2026} (2026).}
\end{table}

\section{Conclusões e Encaminhamentos}
A normalização do \textit{dataset} e a inclusão do histórico de bolsas reduziram o erro de previsão do modelo (RMSE) em 12\%. Para a próxima semana, o foco será transpor essa base enriquecida para uma arquitetura de rede neural recorrente (LSTM), visando capturar as dependências de longo prazo entre o ingresso do aluno e a conclusão da titulação.
